% Chapter 7: The Erd\H{o}s Distinct Distances Problem
\let\LocalMacro\relax

\chapter{Erd\H{o}s Distinct Distances Problem}

\section{The Problem}

\subsection{Informal}

\textbf{Background.} When you place several points on a flat surface, you can measure the distance between every pair. Some pairs may share the same distance, while others are all different. A natural question is how few distinct distances can arise, no matter how cleverly you arrange the points.

\textbf{Given.} You are given $n$ points placed anywhere on a flat surface, with no two points at the same location.

\textbf{Goal.} Find the smallest number of distinct distances that must appear among all pairs of points, regardless of how the points are arranged. That is, determine a lower bound that holds for every possible arrangement.

\textbf{Constraints.} The points can be placed anywhere on the plane---in a grid, on a line, or in any irregular pattern. The bound must hold for all $n$, no matter how large.

\textbf{Example.} Place 4 points at the corners of a square. The distances are: the side length (4 pairs) and the diagonal length (2 pairs). So you see exactly 2 distinct distances. With 4 points on a line at positions 0, 1, 3, 6, the pairwise distances are 1, 2, 3, 3, 5, and 6---giving 5 distinct distances. The question asks: as $n$ grows, what is the minimum number of distinct distances you are forced to see?

\subsection{Formal}
Erd\H{o}s proved in 1946 that $n$ points in the plane determine at least $cn/\sqrt{\log n}$ distinct distances for some constant $c > 0$, and he conjectured this was essentially tight. Guth and Katz proved in 2010 that the answer is at least $cn/\log n$---nearly matching the upper bound.

\subsection{Historical Context}
Erd\H{o}s was one of the most prolific mathematicians in history, publishing over 1,500 papers. He offered prizes for open problems and could pose beautiful questions on any topic. The distinct distances problem was one of his favorites.

\subsection{What Was Known}
Erd\H{o}s proved in 1946 that a square grid of $n$ points produces about $n/\sqrt{\log n}$ distinct distances, and he conjectured this was essentially optimal. For decades, the best known upper bound was roughly $n/\log^{1/4} n$, which was far from the conjectured answer. Guth and Katz's 2010 breakthrough used the polynomial method---a tool from algebraic geometry---to close the gap almost completely.

\subsection{Why It Matters}
The distinct distances problem reveals fundamental truths about how points can be arranged in space. It connects to sensor networks (how to place sensors to avoid redundant measurements), coding theory (how to pack information efficiently), and the study of geometric patterns in nature.

\subsection{Simplified Version}
If all $n$ points lie on a straight line, there are only $n-1$ distinct distances (the gaps between consecutive points). The question is how much the two-dimensional plane forces more. Erdős showed that $n$ points in the plane must produce at least $c \cdot n/\sqrt{\log n}$ distinct distances, and Guth and Katz proved in 2010 that the answer is at least $c \cdot n/\log n$---nearly matching the upper bound from the square grid.

\section{Mathematical Theory}


\subsection{Prerequisites}
This chapter assumes familiarity with:
\begin{itemize}
\item Basic combinatorics and counting arguments
\item Elementary geometry (distances in the plane)
\item Polynomial algebra (degree, roots)
\end{itemize}

\subsection{The Grid Construction}

\begin{definition}[Integer Grid]
Place $n = m^2$ points at integer coordinates $(i,j)$ for $1 \leq i,j \leq m$. The squared distances are $a^2 + b^2$ for $0 \leq a,b \leq m-1$. The number of distinct squared distances is at most $2m^2 = 2n$, but many of these are not achievable as sums of two squares. A careful analysis shows the number of distinct distances is approximately $n/\sqrt{\log n}$.
\end{definition}

This construction shows $f(n) \leq O(n/\sqrt{\log n})$. The lower bound was the hard part.

\subsection{Previous Bounds}

\begin{itemize}
\item Erd\H{o}s (1946): the square-grid construction gives $f(n) \leq Cn/\sqrt{\log n}$, while his counting argument proves the matching-order lower bound $f(n) \geq cn/\sqrt{\log n}$ \cite{erdos1946}
\item Chung (1984): $f(n) \geq cn^{2/5}$ \cite{chung1984}
\item Clarkson, Edelsbrunner, Guibas, Stolfi, and Welzl (1990): $f(n) \geq cn^{2/3}/\log n$ \cite{clarkson1990}
\item Solymosi and T\'oth (2003): $f(n) \geq cn^{2/3}/(\log n)^{c}$ \cite{solymosi2003}
\end{itemize}

Progress was incremental and frustratingly slow.

\subsection{The Polynomial Method}

The proof uses incidence geometry and polynomial methods, but not by constructing a polynomial vanishing on the original point set. Guth and Katz recast equal-distance quadruples as incidences among curves in the three-dimensional space of planar rigid motions, then use polynomial partitioning and algebraic control of the resulting curves and surfaces \cite{guth2015}.

\begin{definition}[Polynomial Method]
Given a set of points $P \subset \mathbb{R}^2$, find a nonzero polynomial $p(x,y)$ of minimum degree that vanishes at all points in $P$. The degree of $p$ constrains the geometry of $P$.
\end{definition}

\section{The Solution}

\subsection{The Big Idea}
The main difficulty is to count quadruples of points that determine the same distance. Guth and Katz use the Elekes--Sharir framework: pairs of point-pairs at equal distance are encoded by incidences among curves in the three-dimensional space of planar rigid motions. Polynomial partitioning controls these incidences, while algebraic arguments handle curves concentrated on low-degree surfaces. This gives a near-linear lower bound for the number of distinct distances.

\subsection{What Was Stopping Progress}
Earlier approaches relied on combinatorial incidence counting and obtained lower bounds with exponent below 1. The breakthrough combined incidence geometry with algebraic tools in a new three-dimensional formulation of the distance-counting problem.

\subsection{The Breakthrough}
Guth and Katz used polynomial partitioning as part of their incidence argument, but their proof is not a direct adaptation of Dvir's finite-field Kakeya proof and does not reduce to a joints count. The key setup is the Elekes--Sharir encoding of equal distances as incidences in the group of planar rigid motions.

\subsection{How It Works}

Erd\H{o}s proved in 1946 that any $n$ points determine at least $cn/\sqrt{\log n}$ distinct distances; a square grid gives an upper bound of order $n/\sqrt{\log n}$. Later lower bounds improved the exponent, but remained below 1 until Guth and Katz proved the lower bound $cn/\log n$.

The problem had resisted methods that counted repeated distances directly. Guth and Katz made progress by encoding equal-distance quadruples as incidences among curves in the rigid-motion group, where polynomial partitioning and algebraic structure could control the count.

Dvir's polynomial proof of the finite-field Kakeya conjecture was an important development in polynomial methods, but it was not the key step in Guth and Katz's distinct-distances proof. Their argument uses the Elekes--Sharir encoding and incidence estimates in the rigid-motion parameter space.

The key innovation was the combination of the Elekes--Sharir framework and polynomial partitioning. The partition is applied in the three-dimensional parameter space for rigid motions, and the proof separately controls incidences on the partition surface using algebraic geometry.

The incidence estimate is not the standard joints theorem. It depends on controlling rich families of curves and their concentration on low-degree algebraic surfaces in the rigid-motion parameter space.

The underlying strategy is to assume there are too few distances, which forces many equal-distance quadruples. The rigid-motion reformulation turns those quadruples into many curve incidences. Polynomial partitioning divides the parameter space into manageable cells, while algebraic structure controls the exceptional incidences on the partition surface.

In particular, the proof does not count circles intersecting a partition curve or derive a bound from the number of such intersections. The relevant incidences are among curves representing point-pairs in the rigid-motion parameter space.

The result is $f(n) \ge cn/\log n$, which almost matches the grid construction's upper bound of $O(n/\sqrt{\log n})$. The $\log n$ gap remains, but the polynomial method closed a gap that had been open for 64 years.

Larry Guth and Nets Katz proved in 2010 \cite{guth2015}:

\begin{theorem}[Guth--Katz, 2010 \cite{guth2015}]
For any set of $n$ points in the plane, the number of distinct distances is at least $cn/\log n$ for some constant $c > 0$.
\end{theorem}

This was essentially optimal (matching Erd\H{o}s's grid construction up to the $\log n$ factor).

Their proof combined several ideas:

\begin{enumerate}
\item \textbf{Elekes--Sharir framework}: Encode equal distances as incidences in the group of planar rigid motions.
\item \textbf{Polynomial partitioning}: Partition the three-dimensional parameter space and estimate incidences inside cells and on the partition surface.
\item \textbf{Algebraic geometry}: Control exceptional curve configurations concentrated on low-degree surfaces.
\end{enumerate}

\begin{highlightbox}
The Guth--Katz proof combined a rigid-motion reformulation, polynomial partitioning, and incidence geometry to obtain the sharp exponent for this problem.
\end{highlightbox}

\subsection{Subsequent Developments}

\begin{itemize}
\item \textbf{Improved bounds}: Later work by Dong, Guth, and others improved the bound slightly.
\item \textbf{Higher dimensions}: The polynomial method has been extended to distinct distances in $\mathbb{R}^d$ for $d \geq 3$.
\item \textbf{Other metric spaces}: The polynomial method has been applied to distinct distances in various settings.
\end{itemize}

\subsection{After the Proof}

The problem is nearly settled but not completely. Guth and Katz proved the lower bound $cn / \log n$, while the square grid construction gives an upper bound of $Cn / \sqrt{\log n}$. These differ by only a factor of $\sqrt{\log n}$, which is remarkably close after 64 years of near-stagnation. The main open question is whether this $\sqrt{\log n}$ gap can be closed. Most experts believe the true answer is $\Theta(n / \sqrt{\log n})$ (matching the grid), but proving the lower bound with that extra $\sqrt{\log n}$ factor has resisted all attempts so far.

The polynomial method that Guth and Katz introduced has become a standard tool in combinatorial geometry. It has been extended to distinct distances in higher dimensions $\mathbb{R}^d$, to distinct distances in other metric spaces, and to a wide range of incidence geometry problems. The joints argument---counting points where multiple curves meet---has found applications far beyond the distinct distances problem. These techniques have influenced work on the Kakeya conjecture, the Bochner--Riesz conjecture, and problems in harmonic analysis.

Several related questions remain open. In three and higher dimensions, the problem of distinct distances is only partially understood: the polynomial method gives bounds but the gap between lower and upper bounds is much larger than in the plane. There are also natural variants---such as distinct distances determined by points on a sphere or on a line---where the answers are not fully known. The question of how many distinct distances can be avoided by clever point configurations continues to drive research in discrete geometry.

\section{Impact}

\begin{itemize}
\item \textbf{Polynomial method}: This proof cemented the polynomial method as a major tool in combinatorial geometry.
\item \textbf{Incidence geometry}: The Guth--Katz approach has influenced many problems in incidence geometry.
\item \textbf{Interdisciplinary connections}: The techniques connect algebra, geometry, and combinatorics in new ways.
\end{itemize}

\section{Computational Verification}

The distinct distances problem can be explored computationally by counting distances for various point configurations and comparing with the theoretical bounds.

\subsection{Counting Distinct Distances}

\begin{verbatim}
import math
from itertools import combinations

def distinct_distances(points):
    """Count the number of distinct distances determined by a set of points."""
    distances = set()
    for (x1, y1), (x2, y2) in combinations(points, 2):
        d = math.sqrt((x1 - x2)**2 + (y1 - y2)**2)
        distances.add(round(d, 10))  # round to avoid floating point issues
    return len(distances)

# Square grid: n = m^2 points
def grid_points(m):
    return [(i, j) for i in range(m) for j in range(m)]

# Points on a line
def line_points(n):
    return [(i, 0) for i in range(n)]

# Points on a circle
def circle_points(n):
    return [
        (math.cos(2*math.pi*k/n), math.sin(2*math.pi*k/n))
        for k in range(n)
    ]

print("Configuration          | n   | distinct distances")
print("-" * 50)
for m in [3, 4, 5, 6, 7, 8]:
    pts = grid_points(m)
    n = len(pts)
    dd = distinct_distances(pts)
    print(f"  {m}x{m} grid            | {n:3d} | {dd}")

for n in [5, 10, 20, 50]:
    dd_line = distinct_distances(line_points(n))
    dd_circle = distinct_distances(circle_points(n))
    print(f"  {n} points on line    | {n:3d} | {dd_line}")
    print(f"  {n} points on circle  | {n:3d} | {dd_circle}")
\end{verbatim}

\subsection{Comparing with Erdos's Bound}

\begin{verbatim}
def erdos_lower_bound(n):
    """Erdos's lower bound: c * n / sqrt(log n)"""
    if n <= 1:
        return 0
    return n / math.sqrt(math.log(n))

def guth_katz_lower_bound(n):
    """Guth-Katz lower bound: c * n / log n"""
    if n <= 1:
        return 0
    return n / math.log(n)

print("Grid: n vs actual, Erdos bound, Guth-Katz bound")
for m in range(2, 11):
    pts = grid_points(m)
    n = len(pts)
    dd = distinct_distances(pts)
    eb = erdos_lower_bound(n)
    gk = guth_katz_lower_bound(n)
    print(f"  n={n:3d}: actual={dd:3d}, Erdos>={eb:.1f}, GK>={gk:.1f}")
\end{verbatim}

\subsection{Optimal Configurations}

\begin{verbatim}
# Hexagonal lattice subset
def hex_points(m):
    pts = []
    for i in range(m):
        for j in range(m):
            x = i + 0.5 * (j % 2)
            y = j * math.sqrt(3) / 2
            pts.append((x, y))
    return pts

print("\nHexagonal lattice configurations:")
for m in [3, 4, 5, 6]:
    pts = hex_points(m)
    n = len(pts)
    dd = distinct_distances(pts)
    print(f"  {m}x{m} hex lattice: n={n}, distinct distances={dd}")
\end{verbatim}

\begin{highlightbox}
Students should run these scripts to observe: (1)~the square grid produces the fewest distinct distances, (2)~collinear points produce even fewer but don't contradict the theorem, and (3)~the Guth--Katz bound $cn/\log n$ is nearly tight for the grid.
\end{highlightbox}

\section{Exercises}

\begin{exercise}[Basic]
Show that $n$ collinear points determine at most $n-1$ distinct distances. Why doesn't this contradict the theorem?
\end{exercise}

\begin{exercise}[Intermediate]
Construct a set of $n$ points on a circle and count the number of distinct distances. How does this compare to the grid construction?
\end{exercise}

\begin{exercise}[Challenge]
Explain why the polynomial method works better in finite fields than in $\mathbb{R}$. What new difficulties arise in the real case?
\end{exercise}


\begin{exercise}[Computational]
Write a Python/SageMath script to explore a concept from this chapter. See the appendix for starter code.
\end{exercise}

\section{Timeline}

\begin{tabularx}{\linewidth}{lX}
\toprule
\textbf{Year} & \textbf{Event} \\ \midrule
1946 & Erd\H{o}s poses the distinct distances problem \cite{erdos1946} \\ 
1984 & Chung improves the lower bound to $n^{2/5}$ \cite{chung1984} \\ 
2008 & Dvir proves the finite field Kakeya conjecture using polynomials \cite{dvir2009} \\ 
2010 & Guth and Katz prove $f(n) \geq cn/\log n$ \cite{guth2015} \\ 
2015+ & Further improvements and extensions \\ \bottomrule
\end{tabularx}

\section{Citations}

\begin{enumerate}
\item Guth, L., \& Katz, N. (2015). ``On the Erd\H{o}s distinct distances problem in the plane.'' Annals of Mathematics, 181(1), 155--190.
\item Erd\H{o}s, P. (1946). ``On sets of distances of $n$ points.'' American Mathematical Monthly, 53(5), 248--250.
\item Dvir, Z. (2009). ``On the size of Kakeya sets in finite fields.'' Journal of the AMS, 23(4), 1107--1119.
\item Chung, F. (1984). ``The number of different distances determined by $n$ points in the plane.'' Journal of Combinatorial Theory, Series A, 36(3), 342--354.
\item Clarkson, K. L., Edelsbrunner, H., Guibas, L. J., Stolfi, M., \& Welzl, E. (1990). ``Combinatorial complexity bounds for arrangements of curves and spheres.'' Discrete \& Computational Geometry, 5(2), 99--160.
\item Solymosi, J., \& T\'oth, C. D. (2003). ``Distinct distances in the plane.'' Discrete \& Computational Geometry, 30(3), 433--448.
\end{enumerate}
